\begin{center}
\normalsize
\setlength{\tabcolsep}{3pt}
\begin{tabular}{@{}lcrp{0.59\textwidth}@{}}
\toprule
Game & Cluster & N & LLM-powered summary \\
\midrule
Dictator & 0 & 487 & Most responses balance self-interest against fairness norms, with a strong plurality choosing a 50/50 split to avoid guilt, signal equity, and minimize regret in an anonymous one-shot interaction. Others rationalize keeping majorities (60-90\%) via token amounts (\$5-40) to ease conscience while maximizing personal gain, reflecting varying altruism levels. \\
\addlinespace[0.4em]
Dictator & 1 & 113 & Most responses reason that since this is a one-shot, anonymous Dictator Game with no repercussions, retaliation, or future interactions, rational self-interest dictates keeping the entire \$100 (allocating \$0). A minority instead weigh fairness/social norms, offering token amounts (\$10-\$50) to balance self-interest against equity concerns. \\
\addlinespace[0.4em]
Prisoner's Dilemma & 0 & 153 & Most responses weigh mutual push’s \$400 cooperative payoff against pull’s safer, self-protective guarantee (\$300+ upside to \$700). Many choose push optimistically, hoping for reciprocity and social efficiency, while others prioritize risk-avoidance and dominant-strategy logic, choosing pull to prevent exploitation given the one-shot, anonymous, no-communication setting. \\
\addlinespace[0.4em]
Prisoner's Dilemma & 1 & 447 & Nearly all responses identify pulling as the strictly dominant strategy, yielding higher payoffs than pushing regardless of the opponent's choice (\$700 vs \$400, \$300 vs \$0). Reasoning emphasizes risk-aversion, avoiding the worst-case \$0 outcome, guaranteeing at least \$300, and maximizing self-interest in a one-shot, non-cooperative, anonymous setting without trust or communication. \\
\bottomrule
\end{tabular}
\end{center}
